\section{Related Works} \label{sec:related-works}
Short-term electricity price forecasting has been studied
extensively over the past decade, with approaches ranging from
classical statistical models to recent deep 
learning architectures.
This section surveys literature most relevant to this paper,
organised around three observations, that prior work tends to
focus on:
\begin{enumerate}
	\item Single-model studies with fixed feature sets.
	\item Multiple-model comparisons operating on single feature
	configurations.
	\item Specialised architectures not evaluated on the Danish
	energy market.
\end{enumerate}
The lack of systematic ablation across feature composition, model architecture, 
and forecasting horizon, particularly for the Danish day-ahead market, is the 
gap this paper seeks to cover.
\subsection{Single-architecture studies with feature-rich inputs}
Various studies train a single model architecture, often with a
fixed feature set chosen for that architecture, and report performance on one 
market. Tschora et al. ~\cite{energyPricePredictionMLNeighbors} forecast 
day-ahead prices in several European markets using gradient boosting trees
on a feature set that includes neighbouring-country prices and
applies SHAP (SHapley Additive exPlanations) to identify the
most predictive inputs. Their study illustrates a common pattern
in the literature, where the results of one, or few, architectures,
typically trained with a single, carefully curated feature set,
are reported as a benchmark rather than as one point within a
sensitivity space. Ziel and Steinert ~\cite{forecastInput} take a related 
approach for mid- and long-term horizons, building a probabilistic model
that connects prices to fundamentals such as temperature,
wind, solar production, load, and renewable penetration. While
these studies establish that market-fundamental features carry
useful information for price prediction, they either primarily
consider the price fluctuations for neighbouring countries, or
do not at all isolate the marginal contribution of each feature
group for price prediction. Laitsos et al. 
~\cite{energyPricePredictionNovelDeepLearning} represent the
deep learning strand of this literature, applying 
attention-based
and neural network-ensemble architectures to short-term price
forecasting. The reported results are strong exhibiting \gls{mape} values in 
the 6\% range,
but the design space is narrow, limiting themselves to only
examining neural network-based architectures on a fixed feature
configuration. The implicit question for a practitioner reading
such papers is whether the architectural sophistication, the
feature engineering, or some interaction of the two drives the
reported accuracy. None of these single or few-model studies
attempt to disentangle the two.

\subsection{Cross-architecture comparisons on fixed feature sets}
A complementary strand of work compares multiple models
on a common dataset. Mishra et al. ~\cite{energyPricePredictionUK} compare 
ARIMA, Prophet, XGBoost, CNN, RNN, and LSTM on UK half-hourly
day-ahead prices using primarily market variables, including
gas prices and energy demand. This is closer in spirit to
this paper in that several model families are evaluated under
identical conditions, but the feature set itself is held constant
across all models. The reader learns which architecture is best
for that particular feature configuration but not whether that
configuration is optimal, whether a simpler one would suffice,
or whether different architectures prefer different feature sets.
These are all questions that this paper addresses directly through
its ablation methodology. El-Azab et al. ~\cite{consumptionPatterns} take a 
step closer to ablation methodology by comparing forecasts with and without
external factors such as temperature and calendar variables,
finding that such inputs improve accuracy. However, the
analysis remains binary, only investigating with versus without
external factors, rather than a structured progression across
feature groups. The granularity of this paper's 13-experiment
progression, applied across eight forecasting horizons, is
intended to fill exactly this gap. For a broader view of the
model-comparison literature, Forootan et al. ~\cite{energySystemsReview} 
provide a review of machine learning and deep learning applications in
energy systems, surveying architectural choices across multiple
forecasting tasks. Such reviews tend to confirm the pattern
observed in the individual papers: most studies specialise in
optimisations within a chosen architecture.

\subsection{Specialised and hybrid architectures}
A separate strand of the literature focuses on hybrid and
decomposition-based methods, in which the input signal is
transformed before being passed to a learning algorithm. Khan
et al. ~\cite{eemdElmPricePrediction} combine Empirical Ensemble Mode 
Decomposition with an Extreme Learning Machine for short-term price
forecasting, exemplifying a class of methods that prioritises
signal-processing sophistication over feature engineering. These
approaches tend to report strong accuracy on specific markets
but introduce methodological complexity that may be difficult
to justify operationally when simpler models with appropriate
feature selection may perform comparably. The present paper
deliberately restricts its scope to general-purpose, off-the-shelf model 
architectures specifically to characterise what is
achievable without such specialisation, providing a baseline
against which more elaborate architectures could be compared
in future work.

\subsection{The gap addressed by this paper}
The literature surveyed above shares three recurring limitations when viewed 
through the lens of practical deployment.
First, results are typically reported for a single feature configuration, 
leaving the question of which features actually matter
unanswered. Second, comparisons between model families,
where they exist, do not also vary feature availability, so it is
unclear whether the relative ranking of architectures is robust to
changes in input data. Third, to the best of our knowledge, no
published study has systematically varied feature composition,
model architecture, and forecasting horizon simultaneously on
the Danish day-ahead market. This is a notable gap given the
market's structural profile, which combines several features
known to complicate price prediction: a small grid with
heavy wind dependence, division into two bidding zones,
and significant cross-border electricity exchange. The present
paper addresses this gap through a 1,052-model ablation
study evaluating eight model architectures across 13 feature
constellations and eight horizons, with two prediction targets,
under default configurations that isolate feature and architecture
effects from hyperparameter tuning.
